\section{Implementação do Protótipo e Coleta de Dados}

Com a formulação completa em programação dinâmica estabelecida, a etapa
seguinte da pesquisa passou a consistir na construção de um protótipo
executável que permitisse confrontar o modelo com instâncias configuradas por
usuários em contexto operacional. Essa etapa não substitui a análise teórica
anterior; ela a complementa. O protótipo funciona como instrumento de
transferência do modelo para uma interface de uso direto e, ao mesmo tempo,
como mecanismo de coleta estruturada de evidências empíricas sobre os tipos de
instância efetivamente testados e sobre o comportamento computacional do
algoritmo.

O protótipo foi implementado como uma aplicação estática em navegador,
publicada via GitHub Pages. A interface recebe uma sequência ordenada de cargas
e uma frota heterogênea definida por classes de veículo. Para cada carga, o
usuário informa os atributos operacionais relevantes; para cada classe de
veículo, informa as capacidades e parâmetros associados ao modelo. A aplicação
executa a formulação completa em programação dinâmica, reconstruindo a
partição ótima da sequência em blocos contíguos e exibindo a alocação obtida.

Do ponto de vista de modelagem, o protótipo corresponde diretamente à
formulação sequencial completa desenvolvida na seção anterior. Em particular, a
implementação considera: (i) a contiguidade dos blocos de carga; (ii) a
heterogeneidade da frota; (iii) o controle explícito do consumo de classes de
veículo quando essa restrição está ativa; (iv) limites de peso; (v) limites de
comprimento com espaçamento entre unidades; (vi) limites máximos de quantidade
de unidades por classe; e (vii) a presença de fretes mínimos faturáveis. A
interface também incorpora regras de validação para impedir configurações
parciais ou semanticamente ambíguas de dados, de modo que a instância enviada
ao algoritmo seja coerente com a estrutura do modelo.

O protótipo oferece duas funções objetivo de interesse prático. A primeira
minimiza o custo total faturado, conforme a estrutura de cobrança definida pelo
modelo. A segunda minimiza o número total de viagens, servindo como variação
operacional relevante para comparação de soluções. Essa duplicidade não altera
a estrutura central do problema, mas amplia o potencial de observação do
protótipo em testes de campo, pois permite identificar em que situações a
preferência por custo e a preferência por quantidade de veículos conduzem a
partições distintas.

Paralelamente ao uso interativo, a aplicação passou a registrar cada execução
de otimização em uma base de dados associada ao frontend. O objetivo dessa
coleta não é auditar usuários, mas preservar instâncias resolvidas e permitir
análise posterior do material produzido nos testes. Cada registro armazena, em
estrutura apropriada para posterior consulta, o objetivo selecionado, os dados
de entrada das cargas, os dados de entrada das classes de veículo, a
configuração efetiva do modelo ativada naquela execução, a classificação da
instância como viável ou inviável, a solução reconstruída quando existente, o
valor da função objetivo e medidas resumidas do esforço computacional, como o
número de estados explorados.

Essa base viabiliza uma etapa empírica posterior da pesquisa. Com ela, torna-se
possível descrever os perfis de instância efetivamente observados nos testes,
identificar quais restrições tendem a ser mais frequentemente ativadas, medir a
incidência de infeasibilidade, comparar a frequência de uso de cada objetivo e
avaliar, ainda que de forma exploratória, como o tamanho das instâncias e a
ativação de diferentes restrições se relacionam com o esforço computacional da
programação dinâmica. Em outras palavras, a coleta transforma o protótipo em
um ambiente de observação sistemática do modelo em uso.

Com isso, a pesquisa passa a ficar organizada em dois eixos complementares. O
primeiro eixo, já desenvolvido neste texto, é teórico-modelador: ele delimita o
problema, posiciona-o na literatura e constrói as formulações relevantes. O
segundo eixo, em andamento, é empírico-exploratório: ele combina a aplicação
das strings de busca ainda pendentes com os testes vivos do protótipo e com a
análise da base de execuções coletadas. A etapa final do trabalho poderá,
portanto, articular a discussão bibliográfica expandida com os resultados
observados no uso real do demonstrador, fechando a conexão entre formulação,
implementação e evidências empíricas.
